\begin{pastexam}[label={px:nsysu-112-dmalg-08}]{中山資工所資安組 112 學年度 第 8 題}
（設計、分析）

(Algorithm points) [10\%] Please describe the algorithm of Merge Sort and analyze its time complexity in
detail.
\end{pastexam}

\begin{pastsoln}
\textbf{演算法}（資結書第~14.4~節）。 合併排序是分治法：把 $A[p..r]$ 從中間切成兩半，兩半各自遞迴排好，再把
兩個排好的序列合併。
\begin{pseudo}{合併排序}
\Procedure{Merge-Sort}{$A, p, r$}
    \If{$p < r$}
        \State $q \gets \lfloor (p + r)/2 \rfloor$
        \State \Call{Merge-Sort}{$A, p, q$}
        \State \Call{Merge-Sort}{$A, q + 1, r$}
        \State \Call{Merge}{$A, p, q, r$}
    \EndIf
\EndProcedure
\Procedure{Merge}{$A, p, q, r$}
    \State 把 $A[p..q]$ 複製到 $L$、$A[q+1..r]$ 複製到 $R$
    \State $i \gets 1$；$j \gets 1$
    \For{$k \gets p$ \To $r$}
        \If{$j > |R|$ \OR ($i \le |L|$ \AND $L[i] \le R[j]$)}
            \State $A[k] \gets L[i]$；$i \gets i + 1$
        \Else
            \State $A[k] \gets R[j]$；$j \gets j + 1$
        \EndIf
    \EndFor
\EndProcedure
\end{pseudo}
呼叫 \proc{Merge-Sort}$(A, 1, n)$。 \proc{Merge} 每一步從兩邊目前最前面的兩個數挑小的放進去。 迴圈不變量：
每一輪開始時，$A[p..k-1]$ 是 $L$、$R$ 裡最小的 $k - p$ 個數，由小到大排好，而 $L[i]$、$R[j]$ 是兩邊剩下的
最小值；所以結束時 $A[p..r]$ 排好了。 相等時取左邊（$\le$），合併排序是穩定的。

例：$\langle 5, 2, 4, 7, 1, 3, 2, 6\rangle$ 切到單一元素後，往上合併成 $\langle 2, 5\rangle\langle 4, 7\rangle
\langle 1, 3\rangle\langle 2, 6\rangle$，再成 $\langle 2, 4, 5, 7\rangle\langle 1, 2, 3, 6\rangle$，最後
$\langle 1, 2, 2, 3, 4, 5, 6, 7\rangle$。

\textbf{時間}。 \proc{Merge} 合併共 $m$ 個數時，迴圈跑 $m$ 次、複製 $m$ 個，是 $\Theta(m)$。 切開只要 $O(1)$。
所以（$n$ 是 $2$ 的次方時，不是的話用上下取整，結果相同）
\[
T(n) = \begin{cases} \Theta(1), & n = 1, \\ 2T(n/2) + cn, & n > 1. \end{cases}
\]
遞迴樹：第 $i$ 層有 $2^i$ 個子問題，每個大小 $n/2^i$，這一層合併的總成本 $2^i \cdot c\,n/2^i = cn$；共
$\lg n + 1$ 層（最下層是 $n$ 個大小 $1$ 的子問題，成本 $\Theta(n)$）。 總和
\[
T(n) = cn \lg n + \Theta(n) = \Theta(n \log n).
\]
也可以用主定理（資結書定理~2.4）：$a = 2$、$b = 2$，$n^{\log_2 2} = n$，$f(n) = \Theta(n)$，第二種情形，
$T(n) = \Theta(n \log n)$。 這個界和輸入的排列無關，最好、平均、最壞情況都是 $\Theta(n \log n)$，達到比較排序
的下界（資結書第~15.1~節）。 合併需要額外的 $\Theta(n)$ 空間，所以合併排序不是原地排序。
\end{pastsoln}
